\section{Discussion and Limitations}
\label{sec:discussion}

The revised flow is deliberately conservative in some places and intentionally incomplete in
others. Those limitations are not incidental; they define where future work should focus.

\subsection{What the Current Design Gets Right}

The strongest engineering choice in the branch is the separation between the behaviours of
OSEK/VDX scheduler and the placement of interruption-related instrumentation. The task
execution order is determined during verification, while the candidate program points are
determined by an explicit analysis step and recorded in a manifest. This keeps the method
maintainable. A change to OSEK/VDX scheduling policy does not automatically require a change
to the instrumentation step, and a change to the candidate-line heuristics does not
automatically require a change to the verification core.

The evidence in \Cref{sec:evaluation} also suggests that the extra structure is not just
architecturally neat; it matters to outcomes. The benchmark report contains multiple cases in
which the improved pipeline is audited as the more faithful model even though it is slower.

\subsection{Current Analysis Limits}

\begin{itemize}[leftmargin=1.5em]
  \item \textbf{Source-file-driven target discovery.} Interleaving functions are selected from
  an explicit source-file list. If the relevant interruption file is not part of the build, the
  analysis cannot discover it.
  \item \textbf{Global-write-centric conflict rule.} Candidate lines are derived from globals
  written by the interleaving function and read or written by the main verification path. This keeps
  the rule simple, but it ignores purely read-side interactions by the interleaving function.
  \item \textbf{\texttt{.c}-only candidate filtering under project scope.} When project-root
  filtering is enabled, candidate recording is restricted to \texttt{.c} files under the root.
  Generated or indirect artifacts outside that slice are ignored.
  \item \textbf{Hand-written manifest IO.} Both the manifest writer and the \texttt{aib}
  reader are hand-coded. There is no schema version field, which makes coordinated changes
  brittle.
\end{itemize}

\subsection{Current OSEK Model Limits}

\begin{itemize}[leftmargin=1.5em]
  \item \textbf{Partial OIL subset.} The parser currently handles the fields needed by the
  implemented scheduling rules: \texttt{TASK}, \texttt{PRIORITY}, \texttt{SCHEDULE},
  \texttt{AUTOSTART}, \texttt{TYPE}, and \texttt{EVENT\_MASK}.
  \item \textbf{AUTOSTART is metadata, not boot behavior.} The configuration records
  \texttt{AUTOSTART}, but the current flow still expects explicit task activation from the
  harness or runtime stubs.
  \item \textbf{\texttt{WaitEvent} is an approximation.} The implementation models the block /
  no-block decision and wake-up dispatch, but not every possible continuation detail of the
  full operating-system behaviour.
  \item \textbf{API coverage is intentionally narrow.} The current branch focuses on task and
  event services relevant to the public benchmarks. Resource, alarm, and interruption-related
  APIs beyond this subset are not yet modeled through the same boundary.
\end{itemize}

\subsection{Performance Trade-offs}

The benchmark report consistently shows that the revised flow pays an overhead tax. Even
public cases add a manifest phase, an injection phase, and a rebuild. Larger cases can still
be worth that cost when the stock baseline either misses a real bug or invents one through
unnecessary interleavings, but the branch is not yet an optimization pass.

This suggests a practical reading of the current status: the branch is already useful as a
more accurate bounded model checking workflow, but not yet as a drop-in faster replacement for
stock CBMC. Improving analysis precision without paying so much orchestration cost is a
natural next step.
